\usepackage[a4paper, total={6in, 8in}]{geometry}
\usepackage{xcolor}
\usepackage{titlesec}
\titleformat{\section}[block]{\color{black}\Large\bfseries\filcenter}{}{1em}{}
\usepackage{fontspec}
\usepackage{hyperref}
\usepackage[normalem]{ulem}
\usepackage{soul}
\usepackage{graphicx}
\usepackage{fancyhdr}
\usepackage{tcolorbox}
\usepackage{luacode}
\usepackage{accsupp}
\usepackage{enumitem}
\usepackage{graphicx} 
\setlist[itemize]{leftmargin=2em}

% Language-specific setup command
\newcommand{\usePreambule}{%
  \ifdefined\useLanguage
  \ifthenelse{\equal{\useLanguage}{PL}}{%
    \usepackage{polski}%
  }{%
    \usepackage[english]{babel}%
  }%
  \setupLanguageSpecifics%
  \else
  \PackageError{preambule}{Language not set}{You must define
  \string\useLanguage\space before calling \string\usePreambule}%
  \fi
}

\usepackage{ifthen}

% Setup language-specific content
\newcommand{\setupLanguageSpecifics}{%
  \ifthenelse{\equal{\useLanguage}{PL}}{%
    % Polish-specific settings
    \newcommand{\exerciseWord}{Zadanie}%
    \newcommand{\exercisesWord}{Zadania}%
    \newcommand{\monthyearcmd}{\protect\directlua{month_year_nominative_pl()}}%
  }{%
    % English-specific settings
    \newcommand{\exerciseWord}{Exercise}%
    \newcommand{\exercisesWord}{Exercises}%
    \newcommand{\monthyearcmd}{\directlua{month_year()}}%
    \hypersetup{pdflang={en}}%
  }%
  \setupFooter%
}

% Setup fancyhdr for all pages
\newcommand{\setupFooter}{%
  \pagestyle{fancy}%
  \fancyhf{}%
  \fancyfoot[C]{\date{\monthyearcmd} - \pdfTitle - [\thepage]}%
  \fancypagestyle{plain}{%
    \fancyfoot[C]{\date{\monthyearcmd} - \pdfTitle - [\thepage]}%
  }%
}

\hypersetup{
  colorlinks=true,
  linkcolor=blue,
  filecolor=magenta,
  urlcolor=blue,
  pdftitle={\pdfTitle},
  pdfpagemode=FullScreen,
}

\colorlet{punct}{red!60!black}
\definecolor{background}{HTML}{EEEEEE}
\definecolor{delim}{RGB}{20,105,176}
\colorlet{numb}{magenta!60!black}


\newtcolorbox{remarkbox}[1][]{
  colback=yellow!10,
  colframe=red!70!black,
  title=Definicja,
  fonttitle=\bfseries,
  boxrule=1pt,
  arc=4pt,
  leftrule=2pt,
  #1
}
\newtcolorbox{infobox}[1][]{
  colback=yellow!10,
  colframe=red!70!black,
  fonttitle=\huge\bfseries,
  fontupper=\Huge\color{red},
  boxrule=1pt,
  arc=4pt,
  leftrule=2pt,
  top=10mm,
  bottom=10mm,
  left=10mm,
  right=10mm,
  center title,
  halign=center,
  valign=center,
  height=\textheight,
  width=\textwidth,
  before skip=0pt,
  after skip=0pt,
  #1
}
\graphicspath{ {config/tex/img/} }
\newcounter{zadanieCounter}

\stepcounter{zadanieCounter}
\newcommand{\Z}{\exerciseWord{} \thezadanieCounter \stepcounter{zadanieCounter}}
\newcommand{\sectionZadanie}{
\section*{\Z}}
\newcommand{\sectionZadania}{
\section*{\exercisesWord}}
\definecolor{codegreen}{rgb}{0,0.6,0}
\definecolor{codegray}{rgb}{0.5,0.5,0.5}
\definecolor{codepurple}{rgb}{0.58,0,0.82}
\definecolor{backcolour}{rgb}{0.95,0.95,0.92}



%opening
\title{\pdfTitle}
\usepackage[cachedir=build]{minted}
\setminted[javascript]{
	frame=lines,
	framesep=2mm,
	baselinestretch=1.2,
	bgcolor=backcolour,
	fontsize=\footnotesize,
	linenos,
	xleftmargin=\parindent
}


\tcbuselibrary{minted}
\newtcblisting{CodeBox}[2][]{
  listing engine=minted,
  minted language=#2,
  minted options={
    linenos,
    breaklines,
    fontsize=\Large
  },
  colback=black!3,
  colframe=black!70,
  title=#1,
  sharp corners,
  boxrule=0.8pt,
  left=6pt,
  right=6pt,
  top=6pt,
  bottom=6pt
}
\newtcblisting{CodeBoxJsx}[1][]{
  listing engine=minted,
  minted language=jsx,
  minted options={
    linenos,
    breaklines,
    fontsize=\Large
  },
  colback=black!3,
  colframe=black!70,
  title=#1,
  sharp corners,
  boxrule=0.8pt,
  left=6pt,
  right=6pt,
  top=6pt,
  bottom=6pt
}
\definecolor{colorLazoYellow1Light}{RGB}{255,248,230}
\definecolor{colorLazoBlue1Light}{RGB}{227,243,255}
\newtcblisting{jscode}{colback=colorLazoBlue1Light,listing
  only,listing engine=minted,minted language=javascript, minted options={
    linenos,
    breaklines
  },
  sharp corners,
  boxrule=0.8pt,
  left=6pt,
  right=6pt,
  top=6pt,
  bottom=6pt
}
\newtcblisting{jsxcode}{colback=colorLazoBlue1Light,listing
  only,listing engine=minted,minted language=jsx, minted options={
    linenos,
    breaklines
  },
  sharp corners,
  boxrule=0.8pt,
  left=6pt,
  right=6pt,
  top=6pt,
  bottom=6pt
}



\newcommand{\MyCopyrightYear}{2024-\the\year}

\begin{luacode*}
  local lfs = require("lfs")

  function include_all_tex_from_folder(folder)
  local files = {}
  for file in lfs.dir(folder) do
  if file:match("%.tex$") then
    table.insert(files, file)
    end
    end
    table.sort(files)

    for _, file in ipairs(files) do
    local path = folder:gsub("\\","/") .. file
    tex.print("\\input{" .. path .. "}")
    end
    end
  \end{luacode*}

  \begin{luacode}
    function month_year_nominative_pl()
    local months = {
      "Styczeń","Luty","Marzec","Kwiecień","Maj","Czerwiec",
      "Lipiec","Sierpień","Wrzesień","Październik","Listopad","Grudzień"
    }
    local t = os.date("*t")
    tex.print(months[t.month] .. " " .. t.year)
    end

    function month_year()
    tex.print(os.date("%B %Y"))
        end
      \end{luacode}

      \newcommand{\IncludeAllTexFromFolder}[1]{%
        \directlua{include_all_tex_from_folder("#1")}%
      }
